\documentclass[12pt]{article}
\usepackage{amsmath,amssymb}
\newcommand{\Dr}{\operatorname{Dr}}
\newcommand{\Trop}{\operatorname{Trop}}
\newcommand{\rk}{\operatorname{rk}}
\newcommand{\cl}{\operatorname{cl}}
\newcommand{\cL}{\mathcal L}
\begin{document}
\section*{Research notes}

Problem: Given a matroid $M$ and the relative Dressian $\Dr(M)$ of
tropical linear subspaces contained in $\Trop(M)$, suppose
$\cL(M)$ has an order-reversing involution $F\mapsto F^\perp$.
The embedded copy of $\cL(M)^{\rm op}$ sends a flat $F$ to the quotient
matroid $M/F\oplus U_{0,F}$.  We seek an order-reversing involution
of $\Dr(M)$ extending this flat opposition.

\subsection*{Conventions}

Let $\mu_M$ be the trivial rank $r=\rk M$ valuated matroid supported
on the bases of $M$.  A rank $k$ point of $\Dr(M)$ is represented by
a rank $k$ valuated matroid $\nu$ with $\nu\preceq\mu_M$, where
$p\preceq q$ means that $p$ is a valuated quotient of $q$.  Inclusion
of projective tropical linear spaces is the opposite order:
\[
\Trop(\theta)\subseteq\Trop(\nu)\Longleftrightarrow \theta\preceq\nu.
\]
Valuated matroids are considered projectively, i.e. up to adding a
common scalar to all finite Plucker coordinates.

For valuated matroids $p,q$ of ranks $a\le b$, the full incidence
relation is
\[
\mathcal I_{a,b}(p,q;S,T):\quad
\min_{i\in T\setminus S}\{p(S\cup i)+q(T\setminus i)\}
\text{ is attained at least twice},
\]
where $|S|=a-1$, $|T|=b+1$, and $S\subseteq T$.  For consecutive
ranks this becomes the three-term relation
\[
 \min\{p(Sx)+q(Syz),p(Sy)+q(Sxz),p(Sz)+q(Sxy)\}.
\]
BEZ Proposition 4.2.3 identifies the incidence-Plucker relations with
valuated quotients, and BEZ Theorem 4.3.1 identifies quotients with
inclusions of projective tropical linear spaces.

Complete flags: for a complete sequence with nonempty layers and
fixed rank $0$ and rank $r$ endpoints, the adjacent three-term
relations alone are the local $M^\natural$-concavity conditions; by
Murota's independent-set cryptomorphism the homogeneous layers are
valuated matroids and adjacent layers are quotients.  Jarra--Lorscheid
Theorem 2.17/Proposition 2.21 gives the flag-matroid-over-the-tropical
hyperfield interpretation.  Joswig--Loho--Luber--Olarte Theorem 3.10
is useful in the finite uniform adjacent-rank case: incidence implies
Plucker.

\subsection*{Structural consequence of the flat-lattice involution}

A finite flat lattice is ranked and semimodular.  If it has an
anti-automorphism $x\mapsto x^\perp$, ranks are reversed:
$\rk(x^\perp)=r-\rk(x)$.  Applying the anti-automorphism to
semimodularity gives the reverse inequality, hence
\[
\rk(x)+\rk(y)=\rk(x\wedge y)+\rk(x\vee y)
\]
for all $x,y$.  Therefore $\cL(M)$ is modular and
\[
(x\vee y)^\perp=x^\perp\wedge y^\perp,\qquad
(x\wedge y)^\perp=x^\perp\vee y^\perp.
\]
The proofs do not need a simplification reduction: loops are in
$\cl(\emptyset)$ and give dependent sets in the coordinate formula;
parallel representatives have equal values by flat constancy.

\subsection*{Flat constancy}

If $\nu\preceq\mu_M$ has rank $k$, then finite $\nu$-bases are
independent in $M$ because the underlying matroid of $\nu$ is a
quotient of $M$.  Also $\nu(B)$ depends only on $\cl_M(B)$ for
independent $k$-sets $B$.

Adjacent-basis proof for $0<k\le r$: compare
$B=A b$ and $B'=A b'$ with $\cl(A b)=\cl(A b')$.  Choose
$C$ of size $r-k$ such that $A b C$ is a basis of $M$; then
$A b' C$ is also a basis.  In
$\mathcal I_{k,r}(\nu,\mu_M;A,A b b' C)$ only the $b,b'$ terms have
finite $\mu_M$ factor, since deleting an element of $C$ leaves the
dependent subset $A b b'$.  Thus $\nu(A b)=\nu(A b')$.  The case
$k=0$ is trivial.

Thus a rank $k$ quotient gives flat coordinates
\[
\bar\nu(X)=\nu(B)\quad\text{for }\rk X=k,\ \cl(B)=X,
\]
and $\nu(A)=\bar\nu(\cl A)$ for independent $A$, while dependent
$A$ has value $\infty$.

\subsection*{One-point completion through a given quotient}

For a single quotient $\nu\preceq\mu_M$ of rank $k$, standard
valuated truncation and relative Higgs elongation insert $\nu$ into a
complete flag:
\[
 h_j(B)=
 \begin{cases}
 \min\{\nu(C): B\subseteq C,\ |C|=k\}, & j\le k,\ |B|=j,\\
 \min\{\nu(C): C\subseteq B,\ |C|=k\}, & j\ge k,\ B\in\mathcal I_j(M),\\
 \infty, & j\ge k,\ B\notin\mathcal I_j(M).
 \end{cases}
\]
This is Murota's valuated truncation/elongation or
$M^\natural$-convex convolution with the indicator of $\mathcal I(M)$.
It gives $h_0\preceq\cdots\preceq h_r=\mu_M$ up to scalar and
$h_k=\nu$.

Important distinction: this one-point completion is not the same as a
general partial-flag completion theorem.  For arbitrary
$\theta\preceq\nu\preceq\mu_M$ of ranks $a<b$, it is not presently
established in these notes that one can find
\[
\theta=h_a\preceq h_{a+1}\preceq\cdots\preceq h_b=\nu .
\]
The earlier broad ``valuated Higgs factorization'' claim should not be
used without proof.  Current literature treats related partial flag
completion questions as nontrivial.

\subsection*{Candidate transform}

For a rank $k$ quotient $\nu$ define
\[
\nu^{\perp_M}(A)=
\begin{cases}
\bar\nu((\cl A)^\perp),& A\text{ independent in }M,\ |A|=r-k,\\
\infty,&\text{otherwise}.
\end{cases}
\]
Adding a scalar to $\nu$ adds the same scalar to $\nu^{\perp_M}$.
The formula is involutive on flat coordinates.

\subsection*{Local elementary opposition}

Let $p\preceq q\preceq\mu_M$ with $\rk p=d$, $\rk q=d+1$.  The local
task is to check the elementary three-term relation for
$q^{\perp_M},p^{\perp_M}$; in a transformed complete flag, the
complete-flag criterion then promotes all adjacent checks to actual
valuated quotients.  If $d=r-1$ the transformed lower rank is $0$ and
there is no adjacent relation to check.  Otherwise fix
$A$ of size $r-d-2$ and $x,y,z\notin A$.  Let
$U=\cl(A)$ and $W=\cl(Axyz)$, with
$\delta=\rk W-\rk U$.

If $A$ is dependent or $\delta\le1$, all relevant terms are $\infty$.

If $\delta=2$, set $G=W^\perp$ (rank $d$) and $H=U^\perp$ (rank
$d+2$).  All finite transformed terms share the same $p$-coordinate
$\bar p(G)$.  A finite term indexed by $e$ has $q$-coordinate
$\bar q(L_e)$ for $L_e=(\cl(Ae))^\perp$, a rank $d+1$ flat in
$[G,H]$.  Terms with $e\in U$ or with the complementary pair not
spanning $W/U$ are $\infty$.  In a rank-two interval either two finite
terms have the same $L_e$, or there are three distinct points.  In the
three-distinct-points case, use $q\preceq\mu_M$: choose a basis
$S$ of $G$, representatives $\ell_x,\ell_y,\ell_z$ in the three
points, and a complement $C$ outside $H$.  The incidence relation for
$q\preceq\mu_M$ with lower set $S$ and upper set
$S\ell_x\ell_y\ell_z C$ has precisely the three finite terms
$\bar q(L_x),\bar q(L_y),\bar q(L_z)$.

If $\delta=3$, set $G=W^\perp$ (rank $d-1$) and $H=U^\perp$ (rank
$d+2$).  Define
\[
P_x=(\cl(Ayz))^\perp,\quad P_y=(\cl(Axz))^\perp,\quad
P_z=(\cl(Axy))^\perp
\]
and
\[
Q_x=(\cl(Ax))^\perp,\quad Q_y=(\cl(Ay))^\perp,\quad
Q_z=(\cl(Az))^\perp.
\]
Modularity in the rank-three interval gives that $P_x,P_y,P_z$ are
distinct atoms over $G$ and
\[
Q_x=P_y\vee P_z,\quad Q_y=P_x\vee P_z,\quad Q_z=P_x\vee P_y.
\]
Choosing representatives $a_x,a_y,a_z$ over a basis $S$ of $G$, the
original three-term relation for $p\preceq q$ is exactly the
transformed relation.

\subsection*{Consequences currently proved}

Using one-point completion of each $\nu\preceq\mu_M$, the elementary
opposition lemma applied along a complete flag through $\nu$ shows
that $\nu^{\perp_M}\preceq\mu_M$.  Thus the formula gives an
involution of the underlying set of $\Dr(M)$.

If $\theta\preceq\nu$ is an elementary quotient, then
$\nu^{\perp_M}\preceq\theta^{\perp_M}$.  Therefore the formula is
order-reversing on every inclusion that admits an elementary
refinement.

For an embedded flat $F$ of rank $f$, $q_F=M/F\oplus U_{0,F}$ has
finite rank $r-f$ flat coordinate at $Z$ iff $Z\vee F=\hat1$.  For an
independent $f$-set $A$ with $Y=\cl(A)$,
\[
\Phi(q_F)(A)<\infty
\Longleftrightarrow Y^\perp\vee F=\hat1
\Longleftrightarrow Y^\perp\wedge F=\hat0
\Longleftrightarrow Y\vee F^\perp=\hat1,
\]
using modularity and complementary ranks.  This is exactly the basis
condition for $q_{F^\perp}=M/F^\perp\oplus U_{0,F^\perp}$.

\subsection*{Main remaining mathematical target}

To answer the original question affirmatively, it remains to prove
one of the following.

(1) Relative elementary-refinement theorem for modular $M$: every
$\theta\preceq\nu\preceq\mu_M$ can be completed to consecutive ranks
between $\theta$ and $\nu$.

(2) Direct full-incidence duality: for arbitrary ranks $a<b$, the
full incidence relations $\mathcal I_{a,b}(\theta,\nu)$ imply the
full incidence relations
$\mathcal I_{r-b,r-a}(\nu^{\perp_M},\theta^{\perp_M})$.

A direct proof should analyze an incidence datum
$A\subseteq T$ for the transformed pair, with
$|A|=r-b-1$, $|T|=r-a+1$, via the modular interval
$[\cl(T)^\perp,\cl(A)^\perp]$.  The nondegenerate case where
$\rk\cl(T)-\rk\cl(A)=b-a+2$ should be the original full incidence
relation after anti-isomorphism.  Degenerate cases of rank
$b-a+1$ appear to require the corank-one quotient relation
$\nu\preceq\mu_M$ localized to a modular interval; this is the main
technical point to formalize.

\subsection*{Computational sanity checks}

Existing checks from previous compute work: full subset-indexed
Plucker and incidence relations were tested for Boolean $U_{4,4}$,
$U_{2,4}$, $U_{2,5}$, Fano $PG(2,2)$ with dot-product polarity,
$U_{2,3}\oplus U_{2,3}$, and random samples in $PG(2,3)$; no
counterexample to flat-coordinate opposition was found.  Additional
useful computations would search in rank $4$ modular examples such as
$PG(3,2)$ for a pair $\theta\preceq\nu$ of rank gap $2$ over a small
value set where transformed full incidence fails.

\end{document}
